\chapter[Análisis de requerimientos]{Especificación: análisis de requerimientos}
\label{cap:especificacion}

\section{Metodología del levantamiento}

Cada requerimiento declara su fuente: \textbf{[PO]} solicitud de la organización destinataria
—«propietario del proyecto» en la terminología de la rúbrica— vía profesor guía; \textbf{[EM]}
estudio de mercado; \textbf{[BT]} base teórica (RFC 8428, ISO/IEC 25010 \parencite{iso25010}, SWEBOK
\parencite{swebok}); \textbf{[DP]} derivación del problema. \textbf{Brecha declarada:} falta la
fuente [ENT], entrevistas directas; su cierre es la acción prioritaria del período.

\section{Síntesis del catálogo}

Se catalogan \textbf{68 requerimientos}: 52 \textbf{funcionales} (RF: lo que el sistema debe
\emph{hacer}) en ocho módulos, y 16 \textbf{no funcionales} (RNF: cómo debe \emph{comportarse} —
rendimiento, seguridad, mantenibilidad) clasificados según
ISO/IEC 25010— con prioridad MoSCoW (49 \emph{Must}, 16 \emph{Should}, 3 \emph{Could}). El catálogo
completo está en el anexo~\ref{anx:requerimientos}.

\begin{table}[H]
\centering\footnotesize
\begin{tabular}{@{}cL{6.2cm}cL{3.2cm}@{}}
\toprule
\textbf{Módulo} & \textbf{Responsabilidad} & \textbf{Req.} & \textbf{Causas} \\
\midrule
M1 & Recepción multiprotocolo (MQTT, HTTP, adaptadores) & 10 & C1, C2 \\
M2 & Normalización al modelo interno (el pipeline) & 7 & C1, C5 \\
M3 & Registro y aprovisionamiento de dispositivos & 6 & C4, C5 \\
M4 & Supervisión de salud (¿quién dejó de reportar?) & 7 & \textbf{C5} \\
M5 & Persistencia y consulta de series de tiempo & 5 & C1, C3, C4 \\
M6 & Alertas y notificación & 4 & C5 \\
M7 & Seguridad, identidad y multi-tenencia & 7 & C1, C3, C4 \\
M8 & Panel de control (la vista, no el todo) & 6 & C1, C2, C4, C5 \\
\bottomrule
\end{tabular}
\caption{Módulos de la solución. El panel es uno de ocho módulos: la sustancia de ingeniería está en
M1--M7.}
\label{tab:modulos}
\end{table}

Requerimientos que anclan el trabajo al dominio de redes: \textbf{RNF-09} (dispositivos en segmento
de red separado), \textbf{RNF-10} (solo conexiones salientes desde los nodos), \textbf{RNF-12}
(añadir una vía de recepción no modifica el núcleo) y \textbf{RNF-16} (la trama nativa cumple SenML).
Y los que sustentan la diferenciación: \textbf{RF-04} (adaptadores por formato), \textbf{RF-24/27/30}
(supervisión de salud) y \textbf{RF-34} (exportación en formato abierto).

\section{Evaluación del mecanismo de transporte}
\label{sec:transporte}

Cinco alternativas sobre seis criterios. Los pesos derivan del problema: la plataforma debe saber
qué llegó y qué no (25\,\%), detectar dispositivos caídos (20\,\%, causa C5), exigir poco al nodo
(20\,\%, lo implementa el compañero), integrarse con poco esfuerzo al backend (15\,\%), pesar poco
por mensaje (10\,\%, la objeción a las cabeceras) y funcionar tras NAT sin abrir puertos (10\,\%,
RNF-10).

\begin{table}[H]
\centering\scriptsize\setlength{\tabcolsep}{2pt}
\begin{tabular}{@{}L{2.7cm}C{1.4cm}C{1.4cm}C{1.3cm}C{1.4cm}C{1.5cm}C{1.1cm}C{1.0cm}@{}}
\toprule
 & \textbf{Entrega confiable} & \textbf{Detecta caídos} & \textbf{Costo en nodo} & \textbf{Facilidad backend} & \textbf{Sobrecarga mensaje} & \textbf{Tras NAT} & \textbf{Total} \\
 & 25\,\% & 20\,\% & 20\,\% & 15\,\% & 10\,\% & 10\,\% & \\
\midrule
\textbf{MQTT} & \textbf{5} & \textbf{5} & 4 & \textbf{5} & \textbf{5} & \textbf{5} & \textbf{4,85} \\
HTTP/REST & 3 & 2 & 4 & \textbf{5} & 2 & \textbf{5} & 3,45 \\
CoAP & 3 & 2 & 4 & 2 & \textbf{5} & 4 & 3,05 \\
WebSocket & 2 & 3 & 3 & 3 & 4 & \textbf{5} & 3,05 \\
Socket TCP propio & 1 & 2 & \textbf{5} & 1 & \textbf{5} & \textbf{5} & 2,60 \\
\bottomrule
\end{tabular}
\caption{Evaluación del mecanismo de transporte (escala 1--5). Los pros y contras en palabras de
cada alternativa están en el anexo~\ref{anx:evaluacion}.}
\label{tab:transporte}
\end{table}

\textbf{Decisión:} MQTT como vía principal —QoS y \emph{Last Will} resuelven de fábrica la entrega
confiable y la detección de caída— y HTTP/REST como vía secundaria para integraciones. El navegador
nunca toca MQTT: consulta la API REST del backend. El socket TCP propio se descarta explícitamente:
lograría la mínima sobrecarga a costa de reinventar sesiones, confirmaciones y formato.

\section{Evaluación del formato de trama}

\begin{table}[H]
\centering\footnotesize
\begin{tabular}{@{}L{2.6cm}L{4.6cm}L{5.6cm}@{}}
\toprule
\textbf{Alternativa} & \textbf{A favor} & \textbf{En contra} \\
\midrule
JSON libre & Flexible, universal & \textbf{Reproduce la causa C1}: cada fuente con sus campos \\
\textbf{SenML (RFC 8428)} & Estándar IETF abierto; identidad, unidades y tiempo resueltos; legible & Algo más extenso que un JSON ad hoc \\
Binario propio & Mínimo peso & Indocumentado, difícil de depurar; solo se justifica en radio \\
CBOR & Compacto y estándar & Ilegible a ojo; librería adicional en el nodo \\
\bottomrule
\end{tabular}
\caption{Evaluación del formato de trama.}
\label{tab:formato}
\end{table}

\textbf{Decisión:} \textbf{SenML como contrato nativo}. Lo que no llegue en SenML entra mediante
\textbf{adaptadores por formato} (uno por formato de fuente, no por sensor); la trama malformada de
un dispositivo conocido va a \textbf{cuarentena} —cruda, con su motivo, sin contaminar las series—;
el dispositivo no registrado se \textbf{rechaza y se registra el intento}. Este es el mismo patrón
de entrada que aplica ThingsBoard con su formato propio \parencite{thingsboardmqtt}; la diferencia
es que aquí el contrato nativo es un estándar abierto.

\section{Arquitectura de recepción y frontera del trabajo}
\label{sec:frontera}

El aporte central es el \textbf{pipeline de tratamiento de tramas} (figura~\ref{fig:pipeline}):
seis filtros independientes y verificables por separado. La identificación de dispositivos combina
el \textbf{registro} (al dar de alta: identidad, tipo, ubicación, credencial) con la
\textbf{autoidentificación} de cada trama (el campo \texttt{bn}); el filtro de validación cruza
ambos.

\begin{figure}[H]
\centering
\resizebox{\linewidth}{!}{%
\begin{tikzpicture}[font=\scriptsize,
  filtro/.style={caja, text width=1.75cm, minimum height=1.5cm, align=center, inner sep=3pt}]

\node[filtro, draw=ambar, line width=1pt] (in) at (0,0) {\textbf{trama}\\[1pt]{\tiny MQTT · HTTP}};

\foreach \i/\nom/\det in {
  1/{autenticar}/{TLS + credencial},
  2/{analizar}/{¿se puede leer?},
  3/{validar}/{contra el registro},
  4/{normalizar}/{al modelo interno},
  5/{supervisar}/{¿quién calla?},
  6/{alertar}/{umbrales y caídas}}{
  \pgfmathsetmacro\x{\i*2.35}
  \node[filtro] (f\i) at (\x,0) {\textbf{\i}\\ \textbf{\nom}\\[1pt]{\tiny \det}};
}
\node[filtro, draw=verdeacento, line width=1.1pt, fill=verdeacento!6] (db) at (16.45,0)
  {\textbf{TimescaleDB}\\[1pt]{\tiny series de tiempo}};

\draw[flechagruesa] (in) -- (f1);
\foreach \i [evaluate=\i as \j using int(\i+1)] in {1,...,5} \draw[flechagruesa] (f\i) -- (f\j);
\draw[flechagruesa] (f6) -- (db);

% rechazos: caída vertical bajo su filtro
\node[notaaviso, text width=2.6cm, align=center, minimum height=1.05cm] (des) at (2.35,-2.3)
  {\textbf{desconocido}\\{\tiny rechazo + registro}};
\node[notaaviso, text width=2.6cm, align=center, minimum height=1.05cm] (cua) at (7.05,-2.3)
  {\textbf{malformada}\\{\tiny cuarentena con motivo}};
\draw[flecha, color=ambar] (f1.south) -- (des.north);
\draw[flecha, color=ambar] (f3.south) -- (cua.north);
\end{tikzpicture}}
\caption{El pipeline de tratamiento de tramas, aporte central del trabajo. Cada filtro es
independiente y verificable por separado; las dos salidas de rechazo preservan la evidencia sin
contaminar las series.}
\label{fig:pipeline}
\end{figure}


El despliegue completo corre como contenedores en un único servidor (figura~\ref{fig:despliegue}),
con reinicio automático ante fallas. \textbf{La frontera con el trabajo del compañero es el tópico
MQTT}: todo lo anterior (sensor, nodo, radio) es de él; todo lo posterior (broker, pipeline, base de
datos, API, panel), de este trabajo. Un \textbf{emulador de nodo} publica tramas en el mismo tópico,
de modo que la validación de la plataforma no depende del hardware ni de terceros.

\begin{figure}[htbp]
\centering
\resizebox{0.98\linewidth}{!}{%
\begin{tikzpicture}[font=\scriptsize]

% lado emisores
\node[caja, text width=2.6cm, align=center, draw=ambar, line width=1pt] (n1) at (0,1.7)
  {sensor + ESP32\\ \tiny (lado del compañero)};
\node[caja, text width=2.6cm, align=center, draw=ambar, line width=1pt] (n2) at (0,0.2)
  {\textbf{emulador de nodo}\\ \tiny (respaldo: publica igual)};
\node[caja, text width=2.6cm, align=center, draw=grislinea] (n3) at (0,-1.3)
  {integración externa\\ \tiny HTTP POST};

% frontera
\draw[line width=1pt, color=gristexto, dashed] (2.3,2.6) -- (2.3,-2.2);
\node[etiqueta, font=\scriptsize\bfseries\color{gristexto}, rotate=90] at (2.05,0.2) {FRONTERA: el tópico};

% servidor
\node[grupo, minimum width=9.6cm, minimum height=5.0cm, anchor=west] (srv) at (3.0,0) {};
\node[font=\scriptsize\bfseries\color{gristexto}, anchor=north west] at (3.1,2.45)
  {UN SERVIDOR (contenedores Docker, \texttt{restart: unless-stopped})};

\node[caja, text width=2.2cm, align=center, draw=verdeacento, line width=1pt] (mq) at (4.6,1.0)
  {\textbf{Mosquitto}\\ \tiny broker MQTT};
\node[caja, text width=2.2cm, align=center, draw=grislinea] (cs) at (4.6,-1.1)
  {ChirpStack\\ \tiny solo si hay rama LoRa};
\node[caja, text width=2.6cm, align=center, draw=azuloscuro, line width=1.1pt, fill=azuloscuro!5] (be) at (8.0,1.0)
  {\textbf{Backend} (NestJS)\\ \tiny pipeline de 6 filtros\\ \tiny + \textbf{API REST}};
\node[caja, text width=2.4cm, align=center, draw=verdeacento, line width=1pt, fill=verdeacento!6] (db) at (11.3,1.0)
  {\textbf{TimescaleDB}\\ \tiny PostgreSQL + series};

\draw[flechagruesa] (n1.east) -- (mq.west|-n1.east);
\draw[flechagruesa] (n2.east) -- (mq.west|-n2.east);
\draw[flecha] (n3.east) -| (8.0,-1.6) -- (be.south);
\draw[flechagruesa] (mq.east) -- (be.west);
\draw[flecha, color=grislinea] (cs.north) -- (mq.south);
\draw[flechagruesa] (be.east) -- (db.west);

% dashboard afuera
\node[caja, text width=2.3cm, align=center, draw=azuloscuro, line width=1pt] (dash) at (8.0,-3.3)
  {\textbf{dashboard}\\ \tiny navegador del usuario};
\draw[flecha] (dash.north) -- node[nota, right=1pt, font=\tiny] {HTTP GET a la API} (be.south|-srv.south);

\end{tikzpicture}}
\caption{Despliegue: todos los componentes corren como contenedores en un único servidor. La
frontera con el trabajo del compañero es el tópico MQTT; el emulador de nodo publica en el mismo
tópico, de modo que la validación no depende del hardware.}
\label{fig:despliegue}
\end{figure}


\section{Evaluación de infraestructura}
\label{sec:infra}

\begin{table}[H]
\centering\footnotesize
\begin{tabular}{@{}L{2.4cm}L{2.4cm}L{7.6cm}@{}}
\toprule
\textbf{Componente} & \textbf{Selección} & \textbf{Razón determinante frente a alternativas} \\
\midrule
Broker MQTT & Mosquitto & C, $\sim$10\,MB de RAM, TLS y autenticación; EMQX/HiveMQ apuntan a millones de conexiones: desproporcionados \\
Persistencia & TimescaleDB & Un solo motor relacional y de series; InfluxDB obligaría a operar dos bases \\
Backend & TypeScript + NestJS & Tipos de la trama definidos una sola vez para backend y panel; módulos 1:1 con M1--M8; OpenAPI automático \\
Rama LoRa & ChirpStack & Autoalojable; publica en MQTT: converge a la vía principal \\
Despliegue & Contenedores Docker & Reproducible con un comando; reinicio automático; la Etapa 4 exige demostrar en condiciones realistas \\
\bottomrule
\end{tabular}
\caption{Decisiones de infraestructura. Las matrices completas están en el repositorio.}
\label{tab:infra}
\end{table}

\textbf{Restricciones:} RE-01 dedicación $\sim$15 h semanales; RE-02 proyecto individual coordinado
con el trabajo del compañero; RE-03 solo componentes abiertos; RE-05 umbral de aprobación del
70\,\%; RE-06 hardware no confirmado, mitigado con el emulador.
